% Imporditud: tools/import_eksperimendid.py
% Allikas: Eesti füüsikaolümpiaadi eksperimendi ülesannete
%   kogu (Rael Kalda, 2023), ülesanne Ü96, lahendus L96.
\setAuthor{EFO žürii}
\setRound{lõppvoor}
\setYear{2006}
\setNumber{G E2}
\setDifficulty{3}
\setTopic{varia}
\setVahendid{vilkuv lamp, stopper}
\setPunktid{12}
\setAllikas{Eksperimendi kogu Ü96, lahendus lk 75}
\setLahendusseis{toores}

\prob{Lamp}
Hinnata valgusallika vilkumise sagedus.

\hint

\solu
Liiguta tuld kiiresti käes edasi tagasi. Püsivalt põleva tule asemel näete punkte. Katsu lugeda punktide hulk ja jätkata võngutamist samas rütmis ning mõõta siis ära võnke periood. Lambi vilkumise sagedus = võngutamise sagedus $\times$ punktide hulk.

Märkus: Täpne lambi vilkumise sagedus on f = 142 Hz.
\probend
